% ============================================================
% BAND XI — DIE ZAHL VOR DEM RECHT · BODY
% ============================================================

\section{Auftrag und Methode}

Dieser Sonderband ist die Faktengrundlage der Podcast-Folge \enquote{Die Zahl vor
dem Recht} (S01-11). Er beantwortet vier Fragen: Was misst der Gender Pay Gap?
Woraus besteht er? Was bleibt nach der Bereinigung --- und was bedeutet dieser
Rest? Und: Wie liest man den internationalen Vergleich, ohne auf die bekannten
Fehlschl\"usse hereinzufallen?

\subsection{Evidenzstufen}

Jede Quelle ist einer von vier Stufen zugeordnet. Verdikte st\"utzen sich
ausschlie\ss lich auf die Stufen A--C; Publikationen politischer Lager (Stufe D)
werden als Argument-Lieferanten zitiert, nie als Beleg.

\begin{longtable}{@{}l L{11.5cm}@{}}
\toprule
\textbf{Stufe} & \textbf{Kriterium} \\
\midrule
A & Amtliche Statistik oder peer-reviewed mit Replikation (Destatis, Eurostat, OECD, ILO; Kleven et al.; Cook et al.) \\
B & Peer-reviewed, Einzelstudie \\
C & Working Paper / graue Literatur mit offengelegten Daten \\
D & Verbands-, Ministeriums- oder Pressepublikation --- nur Wegweiser auf A--C \\
\bottomrule
\end{longtable}

\subsection{Pr\"ufverfahren}

Jede Kernbehauptung durchl\"auft vier Schritte: \textbf{Behauptung} $\rightarrow$
\textbf{Beweislage} (Zahlen mit Quelle) $\rightarrow$ \textbf{Verdikt}
(\vbelegt\ / \vteil\ / \virref\ / \vfalsch\ / \vunent) $\rightarrow$
\textbf{beste Gegenposition}. Abbruchkriterium ist die Ersch\"opfung der aktiven
Widerspruchssuche, nicht die erste Best\"atigung. Die Destatis-Kernbefunde wurden
zus\"atzlich adversarial dreifach gegengepr\"uft (3:0-Voten).

\section{Die zwei Zahlen}

\begin{datenbox}[BEFUND · DESTATIS, VERDIENSTERHEBUNG 2025]
Frauen verdienten 2025 im Durchschnitt \textbf{22,81~Euro} brutto pro Stunde,
M\"anner \textbf{27,05~Euro} --- Differenz 4,24~Euro. Der \textbf{unbereinigte
Gender Pay Gap} betr\"agt damit \textbf{16~\%} (2024: 16~\%, 2023: 18~\%). \belegt
\end{datenbox}

Der unbereinigte Gap vergleicht die Durchschnitts-Stundenl\"ohne \emph{aller}
besch\"aftigten Frauen mit denen \emph{aller} besch\"aftigten M\"anner --- \"uber
alle Berufe, Branchen, Positionen und Arbeitszeitmodelle hinweg. Er misst
ausdr\"ucklich \emph{nicht} die Bezahlung gleicher Arbeit.

\begin{datenbox}[BEFUND · DESTATIS 2025, BEREINIGTER GAP]
Bei vergleichbarer T\"atigkeit, Qualifikation und Erwerbsbiografie verdienten
Frauen 2025 pro Stunde \textbf{6~\%} weniger als M\"anner (westliche
Bundesl\"ander 6~\%, \"ostliche Bundesl\"ander \textbf{9~\%}). Rund 60~\% des
unbereinigten Abstands sind durch beobachtbare Merkmale erkl\"arbar. \belegt
\end{datenbox}

\begin{merksatz}
Destatis deklariert den bereinigten Gap ausdr\"ucklich als \textbf{Obergrenze}
f\"ur m\"ogliche Verdienstdiskriminierung: Weil lohnrelevante Merkmale wie
Erwerbsunterbrechungen in der Erhebung fehlen, fiele der Wert bei deren
Ber\"ucksichtigung \emph{niedriger} aus. Die Sechs ist ein unerkl\"arter Rest ---
kein Diskriminierungsnachweis.
\end{merksatz}

\begin{gegen}
Die Obergrenzen-Deklaration benennt nur die eine Richtung. Methodisch gilt auch
die Gegenrichtung: Eine Bereinigung kontrolliert mit Variablen, die selbst
Diskriminierungs\emph{folgen} sein k\"onnen. Wer wegen zweier Elternzeiten nie
bef\"ordert wurde, dessen fehlende F\"uhrungsposition \enquote{erkl\"art} den
niedrigeren Lohn --- und verdeckt dabei genau den Vorgang, nach dem gefragt wird.
Der bereinigte Gap ist deshalb weder Beweis noch Freispruch. \plausibel
\end{gegen}

\section{Zeitreihe und Ost/West}

Seit Messbeginn 2006 ist der unbereinigte Gap von \textbf{23~\%} (West inkl.
Berlin: 24~\%, Ost: 6~\%) auf \textbf{16~\%} gefallen. \belegt\ Der R\"uckgang
von 18 auf 16~\% zwischen 2023 und 2024 ist teilweise methodisch bedingt
(Umstellung der Verdiensterhebung). \plausibel

Der Ost/West-Kontrast ist das wichtigste Inlands-Exponat dieses Bandes:

\begin{datenbox}[BEFUND · ZWEI KAMMERN DESSELBEN GERICHTS]
Der Osten zeigt den \emph{kleineren unbereinigten} Gap (2006: 6~\% gegen 24~\%
West) --- aber den \emph{gr\"o\ss eren bereinigten} Rest (2025: 9~\% gegen 6~\%
West). Gleiche Rechtslage, andere Erwerbs- und Betreuungsgeschichte: weniger
Funktionspreis, mehr Unerkl\"artes. \belegt
\end{datenbox}

Wer nur die erste Zahl liest, h\"alt den Osten f\"ur gerechter; wer beide liest,
sieht, dass die beiden Gr\"o\ss en Verschiedenes messen. Zugleich zeigt der
Befund: Der \enquote{Preis der Funktion} ist kein Naturereignis --- ein anderer
gesellschaftlicher Bauplan erzeugt einen anderen Preis.

\section{Die Zerlegung: Woraus die L\"ucke besteht}

Rund 63~\% des Abstands 2024 (2,58 von 4,10~Euro) lassen sich statistisch
zuordnen. Die gr\"o\ss ten Einzelposten: \belegt

\begin{longtable}{@{}L{6.2cm}rr@{}}
\toprule
\textbf{Faktor (Destatis-Zerlegung 2024)} & \textbf{Anteil} & \textbf{in Euro/h} \\
\midrule
Berufs- und Branchenwahl & \(\sim\)21~\% & 0,87 \\
Teilzeitquote der Frauen & \(\sim\)19~\% & 0,79 \\
Anforderungsniveau der T\"atigkeit & \(\sim\)12~\% & 0,48 \\
\"ubrige beobachtete Merkmale (u.\,a.\ F\"uhrung, Erwerbsjahre, Zulagen) & Rest bis \(\sim\)63~\% & --- \\
\midrule
\textbf{unerkl\"arter Rest (= bereinigter Gap)} & \(\sim\)37--40~\% & 1,52--1,71 \\
\bottomrule
\end{longtable}

\begin{gegen}
\enquote{Erkl\"art} hei\ss t \emph{zugeordnet}, nicht \emph{gerechtfertigt}. Drei
Anschlussfragen bleiben offen: (1)~Warum zahlen frauentypische Berufe schlechter
--- die Devaluationsforschung findet Hinweise, dass Berufe an Lohn verlieren,
wenn ihr Frauenanteil steigt? \umstritten\ (2)~Wie frei ist die
\enquote{Wahl} der Teilzeit bei l\"uckenhafter Betreuungsinfrastruktur?
(3)~Wessen Lebenslauf tr\"agt die Unterbrechungen --- und warum immer derselbe?
Wer bei der Erkl\"arung stehen bleibt, erkl\"art die Funktionszuweisung zur
Naturkonstante.
\end{gegen}

\section{Der Kinder-Knick (Child Penalty)}

Die sauberste Kausalevidenz zur L\"ucke stammt aus Ereignisstudien um die Geburt
des ersten Kindes: Vor dem Kind laufen die Verdienstkurven von Frauen und
M\"annern nahezu parallel; mit dem Kind knickt die Frauenkurve ein und bleibt
unten, w\"ahrend die M\"annerkurve praktisch unber\"uhrt weiterl\"auft.

\begin{datenbox}[BEFUND · KLEVEN ET AL. 2019 (NBER W25524 / AEA P\&P)]
Langfristiger Einkommensverlust der M\"utter nach dem ersten Kind:
D\"anemark \textbf{21~\%}, Schweden \(\sim\)26~\%, USA/Gro\ss britannien
31--44~\%, \"Osterreich 51~\% --- \textbf{Deutschland 61~\%}, der tiefste Knick
aller untersuchten L\"ander. \belegt
\end{datenbox}

\begin{datenbox}[BEFUND · KLEVEN/LANDAIS/S{\O}GAARD 2019 (AEJ:APPLIED) + CHILD PENALTY ATLAS (RESTUD 2025)]
D\"anemark, Voll-Administrativdaten 1980--2013: Der Anteil der gesamten
Geschlechter-Verdienstungleichheit, der auf Kind-Effekte zur\"uckgeht, stieg von
\(\sim\)40~\% (1980) auf \textbf{\(\sim\)80~\% (2013)} --- fast die gesamte
verbliebene L\"ucke ist kindbedingt. Der \emph{Child Penalty Atlas}
(134 L\"ander, REStud 92(5), 2025) verallgemeinert: Vor der Elternschaft parallele
Erwerbstrends, danach scharfe, dauerhafte Divergenz --- und mit steigendem
Entwicklungsstand werden Kind-Effekte zum \textbf{dominanten Treiber} der
Geschlechterungleichheit. \belegt
\end{datenbox}

Der Befund deckt sich mit dem Destatis-Hinweis, dass Erwerbsunterbrechungen die
wichtigste \emph{fehlende} Variable der Bereinigung sind: Der Kinder-Knick ist
ein zentraler --- in entwickelten L\"andern der dominante --- Treiber der
Gesamtl\"ucke. \belegt\ Ob er der \emph{einzige} Resttreiber ist, bleibt
unentscheidbar: Die Kinder-Strafe kann selbst M\"utter-Diskriminierung enthalten;
die Ereignisstudien dekomponieren das nicht. \vunent

\begin{gegen}
Der Knick selbst ist \"uberwiegend keine Lohnsetzung eines Arbeitgebers ---
er entsteht in Stunden, Unterbrechungen und Jobwechseln. Aber die Ereignisstudien
beantworten nicht, \emph{warum} stets dieselbe Kurve knickt. Pr\"aferenz, Norm
oder Zwang sind in den Registerdaten nicht unterscheidbar. \umstritten
\end{gegen}

\section{Der Rest und seine Ausleuchtung: Feld- und Naturexperimente}

\begin{datenbox}[BEFUND · COOK, DIAMOND, HALL, LIST, OYER (REV. ECON. STUDIES 2021)]
\"Uber eine Million Uber-Fahrer, geschlechtsblinder Verg\"utungsalgorithmus:
dennoch \textbf{\(\sim\)7~\%} Verdienstl\"ucke pro Stunde --- restlos
zur\"uckf\"uhrbar auf drei Faktoren: Plattform-Erfahrung, Ortswahl (Wohnort,
Sicherheitsempfinden), Fahrtempo. Keine Diskriminierung durch Lohnsetzung oder
Kunden nachweisbar. \belegt
\end{datenbox}

Der Uber-Befund schneidet in beide Richtungen: Er zeigt, dass eine L\"ucke ganz
ohne diskriminierende Lohnsetzung entstehen kann --- und dass \enquote{gleicher
Tarif} sie nicht schlie\ss t, solange Lebensumst\"ande ungleich verteilt sind.

Audit- und Korrespondenzstudien (identische Bewerbungen mit und ohne
Elternschafts-Signal) zeigen wiederholt Nachteile f\"ur M\"utter im R\"ucklauf
--- ein Mechanismus, der \emph{vor} jedem Lohnzettel wirkt und in keiner
Lohnzerlegung auftaucht. \plausibel\ F\"ur Deutschland fehlen Feldexperimente in
gro\ss er Breite; ihre Nachlieferung ist Teil des Falsifikationsangebots
(\S\,9).

\begin{datenbox}[BEFUND · BLAU/KAHN 2017 (JEL) UND SOEP-WEST 1984--2020]
Der Standard-Survey des Fachs (Blau \& Kahn, JEL 55(3)) schlie\ss t aus der
Experimentallage: \enquote{discrimination cannot be discounted} --- direkte
Diskriminierung ist als Teil des Rests nicht auszuschlie\ss en. Und die deutsche
Langzeit-Zerlegung (Bonaccolto-T\"opfer/Castagnetti/Rosti 2023, SOEP West
1984 vs.\ 2020, RIF-Quantile) findet: Die L\"ucke sank \"uber fast die gesamte
Verteilung --- \textbf{am oberen Ende stieg sie signifikant}, getrieben nicht von
Merkmalen, sondern von unterschiedlichen \emph{Renditen} f\"ur gleiche Merkmale
(Vollzeiterfahrung, Betriebszugeh\"origkeit): der Gl\"aserne-Decke-Befund. \belegt
\end{datenbox}

\section{Der Weltvergleich --- und wie man ihn liest}

Internationale Gap-Werte sind \emph{nicht} direkt stapelbar: Eurostat misst
unbereinigte Brutto-Stundenl\"ohne, die OECD den Median der
Vollzeitbesch\"aftigten. Jede L\"anderzahl ist daher nur im Dreiklang lesbar:
\textbf{L\"ucke + Frauenerwerbsquote + Rechtsstellung}.

\begin{longtable}{@{}L{3.4cm}rL{2.6cm}L{6.2cm}@{}}
\toprule
\textbf{Land} & \textbf{Gap} & \textbf{Frauen-\newline erwerbsquote}\footnotesize{ (ILO 15+)} & \textbf{Lesehinweis} \\
\midrule
Deutschland & 16~\% (Destatis) / 17,6~\% (Eurostat) & hoch & Teilzeitkultur; bereinigt 6~\% \\
\"Osterreich & 18,3~\% (Eurostat) & hoch & wie DE: Teilzeit-Hochlohnland \\
Belgien & 0,7~\% (Eurostat) & hoch & Tarifbindung, Lohnkompression \\
Luxemburg & $-$0,9~\% (Eurostat) & hoch & Komposition: hochqualifizierte Dienstleistungsjobs \\
EU-Schnitt & 12~\% (Eurostat) & --- & Referenz \\
Italien & 2,2~\% (Eurostat) & 41,3~\% & \textbf{Selektions-Paradox}: kleine gemessene L\"ucke, weil viele Frauen gar nicht erwerbst\"atig sind \\
T\"urkei & niedrig gemessen & 35,8~\% & Selektion wie Italien \\
Mexiko & 16,7~\% (OECD) & 46,8~\% & Selektion + informeller Sektor \\
Brasilien & --- & 53,2~\% & Datenlage heterogen \umstritten \\
USA & \(\sim\)16~\% (OECD) & mittel & Median Vollzeit; andere Messlatte \\
Kanada & 17,1~\% (OECD) & mittel & Pay-Equity-Gesetzgebung als Kontrast \\
Australien & \(\sim\)11~\% (OECD) & mittel & --- \\
Japan & 22,0~\% (OECD) & mittel & formale Rechtsgleichheit, Dauerverf\"ugbarkeits-Norm \\
S\"udkorea & 29,3~\% (OECD) & mittel & gr\"o\ss te L\"ucke der Industriewelt \\
China & --- & hoch & post-sozialistische Erosion; keine belastbare amtliche Reihe \umstritten \\
Russland & --- & hoch & hohe Erwerbstradition; Datenlage eingeschr\"ankt \umstritten \\
Saudi-Arabien & --- & 34,1~\% & Die L\"ucke beginnt vor dem Lohn: Zugang und Recht; Erwerbsquote nach Reformen fast verdoppelt \\
\bottomrule
\end{longtable}

\begin{merksatz}
Ein kleiner gemessener Gap kann R\"uckstand verstecken (Selektion), ein
gro\ss er kann Teilhabe enthalten. Die L\"ucke ist kein Gerechtigkeits-Ranking
--- sie ist ein R\"ontgenbild, und R\"ontgenbilder liest man nicht mit einem
Blick.
\end{merksatz}

\begin{gegen}
Auch die Dreiklang-Lesung hat Grenzen: Erwerbsquoten (ILO, 15+) sind
demografieabh\"angig, Rechtsindizes (World Bank \emph{Women, Business and the
Law}) messen Gesetzestexte, nicht Praxis --- und die Revision des Index (WBL
1.0 auf 2.0) verschiebt Punktwerte erheblich. Dieser Band verzichtet deshalb auf
WBL-Punktangaben und nutzt den Index nur ordinal. \plausibel
\end{gegen}

\section{Die Verdikt-Tafel}

Neun \"offentliche Behauptungen, gepr\"uft im Vier-Schritt-Verfahren
(vollst\"andige Beweislage: \texttt{studie/verdikte-ENTWURF.md} im Repositorium).

\begin{longtable}{@{}L{8.6cm}l@{}}
\toprule
\textbf{Behauptung} & \textbf{Verdikt} \\
\midrule
\enquote{Frauen verdienen f\"ur gleiche Arbeit 21~\% weniger.} & \virref \\
\enquote{Der Gender Pay Gap betr\"agt 21~\%.} (als aktuelle Zahl) & \vfalsch \\
\enquote{Die bereinigten 6~\% sind der bewiesene Diskriminierungsanteil.} & \virref \\
\enquote{(Fast) alles erkl\"arbar --- also keine Diskriminierung, kein Problem.} & \virref \\
\enquote{Der Equal Pay Day markiert, ab wann Frauen umsonst arbeiten.} & \vteil\ / \virref \\
\enquote{Der Kinder-Knick ist ein zentraler Treiber.} & \vbelegt \\
\enquote{Frauen verhandeln schlechter --- selbst schuld.} & \vteil\ / \virref \\
\enquote{In L\"andern mit kleinerem Gap geht es Frauen besser.} & \virref \\
\enquote{Ost/West zeigt: gleiche Rechte, anderer Funktionspreis.} & \vteil \\
\bottomrule
\end{longtable}

Bemerkenswert ist die Symmetrie: Die Plakatformel der einen Seite und die
Entwarnungsformel der anderen tragen \emph{dasselbe} Verdikt --- \virref\ ---
nur spiegelverkehrt angezogen.

\section{Bewertung der Arbeitsthese}

\textbf{These (Auftraggeber):} Deutschland ist bei der Rechtsgleichheit weiter
als die meisten Vergleichsl\"ander; die verbleibende unbereinigte L\"ucke ist
\"uberwiegend der \"okonomische Preis f\"ur die Abweichung von der traditionellen
Funktionsteilung --- nicht direkte Lohndiskriminierung.

\textbf{Belege:} Formale Rechtsgleichheit vollst\"andig (Band X:
Pflicht-Grammatik-Zeiger auf null; WBL-Spitzengruppe). \plausibel\ Bereinigter
Gap 6~\% als Obergrenze. \belegt\ Kinder-Knick 61~\% als tiefster der
Vergleichswelt. \belegt\ Ost/West-, Selektions- und Ostasien-Befund zeigen
dieselbe Handschrift: Nicht das Recht setzt die L\"ucke, die Funktionszuweisung
setzt sie. \plausibel

\textbf{St\"arkste Gegenbelege (gegen die Ausschlie\ss lichkeit):} Blau/Kahn 2017:
Diskriminierung \enquote{cannot be discounted} \belegt; SOEP-West: am oberen Ende
gewachsene L\"ucke mit ungleichen Renditen f\"ur gleiche Merkmale (Gl\"aserne
Decke) \belegt; die Kinder-Strafe ist nicht in Pr\"aferenz vs.\
M\"utter-Diskriminierung dekomponierbar \belegt; Audit-Evidenz zu
M\"utter-Nachteilen vor dem Lohnzettel \plausibel; Devaluation frauentypischer
Berufe \umstritten; der gr\"o\ss ere \emph{bereinigte} Ost-Rest.

\begin{merksatz}
Verdikt: \vteil\ (Konfidenz mittel--hoch) --- das \emph{\"uberwiegend} tr\"agt,
das \emph{nicht} tr\"agt nicht: Der Mechanismus (Funktionspreis) erkl\"art den
Gro\ss teil der L\"ucke, aber direkte Diskriminierung ist nach Experimentallage
und Gl\"aserne-Decke-Befund nicht auszuschlie\ss en. Und der Preis ist kein
Naturereignis: Der Ostbefund zeigt, dass ein anderer Bauplan einen anderen
Preis erzeugt.
\end{merksatz}

\textbf{Falsifikationsangebot} (im Stil von Band X): (1)~Feldexperimente in
deutscher Breite --- br\"ache der M\"utter-Nachteil im R\"ucklauf aus, fiele der
Diskriminierungsanteil der Restl\"ucke. (2)~Lohnbiografien unter symmetrischer
Betreuung, Steuer und Elternzeit --- verschwindet der Rest dann nicht, war es
nie der Funktionspreis; verschwindet er, war es nie das Vorurteil. Beides
w\"are Erkenntnis.

\section{Quellen}

\textbf{Stufe A --- amtlich / peer-reviewed repliziert:}
\begin{itemize}
\item Statistisches Bundesamt: Gender Pay Gap 2025 --- unbereinigt 16~\%, bereinigt 6~\% (West 6~\%, Ost 9~\%), Erkl\"arungsanteil \(\sim\)60~\%; Obergrenzen-Deklaration. Pressemitteilung PD25\_453 (Dez.\ 2025) und Themenseite Verdienste/GenderPayGap. \url{https://www.destatis.de/DE/Presse/Pressemitteilungen/2025/12/PD25_453_621.html}
\item Statistisches Bundesamt: Gender Pay Gap 2024 --- 16~\% (22,24 vs.\ 26,34~Euro), Zerlegung (Beruf/Branche 21~\%, Teilzeit 19~\%, Anforderungsniveau 12~\%). PD25\_056 (Feb.\ 2025). \url{https://www.destatis.de/DE/Presse/Pressemitteilungen/2025/02/PD25_056_621.html}
\item Eurostat: Gender pay gap in unadjusted form, 2023 --- EU 12,0~\%; DE 17,6~\%; AT 18,3~\%; BE 0,7~\%; IT 2,2~\%; LU $-$0,9~\%. \url{https://ec.europa.eu/eurostat/statistics-explained/index.php?title=Gender_pay_gap_statistics}
\item OECD: Gender wage gap (Median, Vollzeit), 2023 --- OECD \(\sim\)11~\%; JP 22,0~\%; KR 29,3~\%; US \(\sim\)16~\%; CA 17,1~\%; MX 16,7~\%; AU \(\sim\)11~\%. \url{https://www.oecd.org/en/data/indicators/gender-wage-gap.html}
\item Kleven, Landais, Posch, Steinhauer, Zweim\"uller (2019): \emph{Child Penalties Across Countries: Evidence and Explanations}. AEA Papers \& Proceedings 109 / NBER w25524 --- Langfrist-Penalty DK 21~\%, DE 61~\%. \url{https://www.nber.org/papers/w25524}
\item Kleven, Landais, S{\o}gaard (2019): \emph{Children and Gender Inequality: Evidence from Denmark}. AEJ: Applied 11(4) --- Kind-Anteil an der Gesamtungleichheit 40~\% (1980) auf 80~\% (2013). \url{https://www.aeaweb.org/articles?id=10.1257/app.20180010}
\item Kleven, Landais, Leite-Mariante (2025): \emph{The Child Penalty Atlas}. Review of Economic Studies 92(5) --- 134 L\"ander; Kind-Effekte als dominanter Treiber in entwickelten L\"andern. \url{https://academic.oup.com/restud/article/92/5/2954/7924096}
\item Blau, Kahn (2017): \emph{The Gender Wage Gap: Extent, Trends, and Explanations}. Journal of Economic Literature 55(3), 789--865 --- \enquote{discrimination cannot be discounted}. \url{https://www.aeaweb.org/articles?id=10.1257/jel.20160995}
\item Cook, Diamond, Hall, List, Oyer (2021): \emph{The Gender Earnings Gap in the Gig Economy}. Review of Economic Studies 88(5), 2210--2238 --- 7~\%-L\"ucke, drei Erkl\"arfaktoren. \url{https://www.restud.com/paper/the-gender-earnings-gap-in-the-gig-economy-evidence-from-over-a-million-rideshare-drivers/}
\item ILO-Modellsch\"atzungen (World Bank Gender Data Portal), Frauenerwerbsquote 15+, 2023: IT 41,3~\%; TR 35,8~\%; MX 46,8~\%; BR 53,2~\%; SA 34,1~\%. \url{https://data.worldbank.org/indicator/SL.TLF.CACT.FE.ZS}
\end{itemize}

\textbf{Stufe B/C --- Einzelstudien / graue Literatur:}
\begin{itemize}
\item Bonaccolto-T\"opfer, Castagnetti, Rosti (2023): RIF-Quantilszerlegung SOEP West 1984 vs.\ 2020, Journal for Labour Market Research 57(11) --- Gl\"aserne-Decke-Befund (oben gewachsener Gap, ungleiche Renditen); Methodik und Grenzen der Bereinigung (deskriptiv; Selektionsverzerrung). \belegt
\item Correll, Benard, Paik (2007): \emph{Getting a Job: Is There a Motherhood Penalty?} --- Audit-Design zu M\"utter-Nachteilen. \plausibel
\item Devaluationsforschung (England; Levanon et al.): Lohnentwicklung bei steigendem Frauenanteil eines Berufs. \umstritten
\item Goldin (2014/2021): \emph{greedy jobs} --- Nichtlinearit\"at der Stundenertr\"age als Mechanismus der Branchenl\"ucken. \plausibel
\end{itemize}

\textbf{Stufe D --- Lager-Publikationen (nur als Argument-Quellen verwendet):}
BMFSFJ, WSI/DGB, Antidiskriminierungsstelle des Bundes; IW K\"oln, INSM; die
sechs Referenzvideos der Folge (Argument-Auswertung:
\texttt{quellen/ABGLEICH.md}).

\vspace{1.5em}
{\color{obsmuted}\rule{\textwidth}{0.4pt}}

{\footnotesize\color{obsmuted} Sonderband XI der Reihe \emph{Die Frau vor dem
Recht / Homo ante legem}. Erstellt als Faktengrundlage der Podcast-Folge
S01-11 \enquote{Die Zahl vor dem Recht}. Kernbefunde adversarial verifiziert
(3:0-Voten); \"ubrige Werte gegen die genannten Prim\"arquellen gepr\"uft.
Lizenz: CC BY-SA 4.0.}
